\section{Functional Verification and Requirements Coverage}
\label{sec:eval-functional}

We verified the functional and non-functional requirements of \cref{ch:requirements} by four means: automated tests, manual end-to-end walkthroughs of the researcher and participant workflows, measurement of the exported data, and inspection of the implemented modules.
\Cref{tab:req-coverage} records the method and verdict for each requirement group.
The backend test suite, which covers session lifecycle, decision-proposal handling, replay recovery, and the export serialisers, passes in full (100 of 100 tests), and that suite also contains the architecture-boundary test discussed in \cref{sec:eval-modularity}.
The frontend has no automated test framework, so the participant- and researcher-facing flows were verified by manual walkthrough rather than by automated tests; we record this openly in the table and revisit it in \cref{sec:eval-threats}.

In \cref{tab:req-coverage} the method codes are T (automated test), W (manual walkthrough), M (measurement of the exported data), I (code inspection), and D (demonstration).
A verdict of ``Met'' means the requirement was verified by one of those means in this work, whereas ``Implemented'' means it is present in the system by inspection but was not exercised by the recorded sessions or the test suite.

\begingroup
\small
\begin{longtable}{p{0.30\linewidth} p{0.42\linewidth} p{0.18\linewidth}}
\caption{Requirements coverage: the method by which each requirement group was verified and the resulting verdict.}
\label{tab:req-coverage} \\
\toprule
\textbf{Requirement group} & \textbf{Method} & \textbf{Verdict} \\
\midrule
\endfirsthead
\multicolumn{3}{c}{\tablename~\thetable{} --- continued} \\
\toprule
\textbf{Requirement group} & \textbf{Method} & \textbf{Verdict} \\
\midrule
\endhead
\midrule
\multicolumn{3}{r}{\footnotesize Continued on next page} \\
\endfoot
\bottomrule
\endlastfoot

\multicolumn{3}{l}{\textit{Functional requirements}} \\
\midrule
FR1 System execution and setup          & Walkthrough of guided setup and readiness gating (W) & Met \\
FR2 Eye-tracker detection and licensing & Real Tobii sessions detected, licensed, streamed (W, D) & Met \\
FR3 Calibration integration             & Calibration metrics recorded and gated on (M, W) & Met \\
FR4 Gaze acquisition and mapping        & Sustained stream and token-level mapping in the record (M) & Met \\
FR5 Reading content and presentation    & Markdown rendering and typographic controls exercised (W) & Met \\
FR6 Researcher interface (second screen) & Live mirror and researcher console operated (W) & Met \\
FR7 Intervention runtime (pluggable)    & Module-contract test and registry inspection (T, I) & Met \\
FR8 Decision-strategy abstraction       & Decision-proposal lifecycle tests and strategy registry (T, I) & Met \\
FR9 Context preservation                & Per-intervention outcomes measured and logged (M, D) & Met \\
FR10 Real-time performance monitoring   & Gaze-validity rate in the telemetry record (M) & Met \\
FR11 Experimental control (modes)       & Manual mode exercised; modes present by inspection (W, I) & Met \\
FR12 Logging and export                 & Sessions reconstructed from the export alone (D, I) & Met \\
FR13 Comprehension quiz                 & Quiz pipeline present; not exercised here (I) & Implemented \\
FR14 Reading-material library           & Library and CRUD present (I) & Implemented \\
FR15 Experiment setup templates         & Template definition present (I) & Implemented \\
FR16 Multi-item session sequencing      & Sequencing and transition events present (I) & Implemented \\
FR17 Sensing-mode configuration         & Tobii and mouse sources both operable (D, I) & Met \\
FR18 Eye-movement analysis layer        & Sessions analysed by an external provider; built-in present (M, D, I) & Met \\
FR19 External module-provider framework & Three providers connect; rejection paths smoke-tested (D, T) & Met \\
FR20 Decision-proposal approval workflow & Proposal-gating lifecycle tests (T, I) & Met \\
FR21 Session replay playback            & Replay reconstructs a session from its record (D, I) & Met \\
FR22 Session-state recovery             & File-backed recovery covered by tests (T) & Met \\
FR23 Per-context reader-shell settings  & Independent per-view settings present (I) & Implemented \\
FR24 Facial-state sensing (optional)    & Optional webcam path present (I) & Implemented \\
\midrule
\multicolumn{3}{l}{\textit{Non-functional requirements}} \\
\midrule
NFR1 Modularity            & Dependency graph and an executable boundary test (T, I) & Met \\
NFR2 Low latency           & Transport and decision-path latencies within budget (M) & Met \\
NFR3 Usability             & Guided setup and unified-console walkthrough (W) & Met \\
NFR4 Reliability           & Provider disconnect logged; session continued (D) & Met \\
NFR5 Reproducibility       & Sessions reconstructed from the record alone (D) & Met \\
NFR6 Extensibility         & Additive module and external-provider integration (I, D) & Met \\
NFR7 Recoverability        & Checkpoint and restore covered by tests (T) & Met \\
NFR8 Hardware independence  & Full operation under mouse simulation (D, I) & Met \\
\end{longtable}
\endgroup

The verdict for NFR2 rests mainly on the client transport latency reported in \cref{sec:eval-performance}; the decision-pipeline and out-of-process decision latencies were exercised in a dedicated advisory-mode validation run, where both also stayed within budget, and we read that evidence as a single short run with a reference provider rather than as a result drawn from the eight behavioural sessions.
